\begin{theorem}[13.1 (формула конечных приращений)]\label{13.1}
% название не устоявшееся
Пусть $E_1, E_2$~--- вещественные банаховы пространства,  $D$~-- выпуклая область, $F: D \to E_2$~--- дифференцируема на $F$.
Тогда $\forall x_0, x_1 \in D$
\[
\|F(x_1) - F(x_0)\| \leqslant \sup_{y \in (x_0, x_1)}
\|F'(y)\| \cdot \|x_1 - x_0\|
\]
\end{theorem}

\begin{proof}
[тут картинка]
$\varphi(t) = f(F(x(t)))$ дифференцируема

\[
\varphi(1) - \varphi(0) = f(F(x_1)) - f(F(x_0)) = 
f(F(x_1) - F(x_0))
\]

\[
    \varphi'(t)=f(F'(x(t))(x_1-x_0))
\]

$f$ не затронута операцией дифференцирования, так как является линейным функционалом.

Воспользуемся следствием 2 из теоремы Хана-Банаха:
$E$~--- нормированное пространство, $0 \neq x \in E$.
Тогда $\exists f \in E^*$, такая что $\|f\| = 1$ и $f(x) = \|x\|$.

Возьмем $f \in E_2^*$ для $(F(x_1)-F(x_0) \in E_2$. Для этого $f$
\[
\varphi(1) - \varphi(0) = \|F(x_1) - F(x_0)\|
\]

По обычной теореме Лагранжа \[ |\varphi'(\xi)(1-0)=|\varphi(1)-\varphi(0)| \]

\[
|f(F'(x_{\xi}) \underbrace{(x_1 - x_0)}_{\Delta x})| \leqslant \underbrace{\|f\|}_{1} \cdot
\sup_{y \in (x_0, x_1)} \|F'(y)\| \cdot \|x_1 - x_0\|
\]

\end{proof}

\begin{exercise}
    Придумать пример $F, G$~--- слабо дифференцируемы, а $H=G \circ F$~--- нет
\end{exercise}

\begin{theorem}[13.2]
    $F'_{\text{сл.}}(x_0)$ существует и непрерывна в окрестности $x_0 \implies \exists F'_{\text{Фр.}}(x_0)=F'_{\text{сл.}}(x_0)$ 
\end{theorem}

\begin{exercise}
$H(\R)$, $F(x) = \|Ax\|^2, A \in \L(H)$. $F' = ?$
\end{exercise}

\begin{statement}[(одно из свойств)]
    Если $F, G$~--- дифференцируемы по Фреше, то $H=G \circ F$~--- дифференцируема по Фреше и $H'(x_0)=...$
\end{statement}

\textbf{Применение теоремы \nameref{13.1}:}
\begin{enumerate}
    \item Дифференциальное исчисления в банаховых пространствах.
    В частности, решение экстремальных задач.
    % книжка Иоффе, Тихомиров, Экстремальные задачи 1974
    Вариационное исчисление: вариация по Лагранжу == слабый дифференциал Гато

    \item Решение уравнений, метод Ньютона,...
    Теорема 13.1~--- удобный способ установления того, что $F$~--- сжимающее отображение.
    В частности: решение уравнений $F(x)=0$ ($x_{n+1}=x_n-[F'(x_n)]^{-1} F(x_n)$. Если предел существует, то $x=x--[F'(x)]^{-1} F(x)$, то есть $F(x)$ должен равняться 0)
    % книжка Канторович, Акилов, Функциональный анализ, Часть II: "Функциональные уравнения"
\end{enumerate}



\subsubsection{Теоремы о неподвижных точках:}
\begin{enumerate}
    \item $\R^n$
    \begin{theorem}[Боля-Брауэра]
        [Картинка] $F$~--- непрерывно и переводит шар в подмножество шара $\implies \exists x \in \overline{B}: Fx=x$
    \end{theorem}
    
\end{enumerate}
% рисуем карту Африки
% кормим римских голубей пиццей

\begin{exercise}
    $l_2$. Придумать непрерывное отображение $F: \overline{B}(0; 1) \to \overline{B}(0;1)$ такое, что у него нет неподвижной точки
\end{exercise}

% Если ты шарик, ты в компактную тарелку не влезешь
% "Здравствуйте, я пришелец" (c) Коновалов
% Выбираем теорему в подарок

\begin{note}
    Чтобы перейти в бесконечно мерное пространство мы должны заплатить большую цену: где-то должно появиться слово компактность. Либо берем непрерывное отображения между компактами (теорема Шаудера), либо компактный оператор на шарах. Второе приятнее, так как шары везде есть
\end{note}

\begin{definition}
$F: E_1 \to E_2$ называется компактным отображением, если  оно непрерывно и для любого ограниченного множества $B$ $F(B)$~--- предкомпакт. Увы, $F$ не обладает непрерывностью.
\end{definition}

\textbf{Примеры применения теоремы Шаудера}

\[
\begin{cases}
y' = f(x,y)\\
y|_{x=0}=y_0
\end{cases}
\]

\begin{theorem}[Фробениус, Перрон]
    $A=(a_{ij})$~-- вещественная матрица $n \times n, a_{ij} > 0 \; \forall i, j$. Тогда у $A$ существует собственный вектор $x=(x_1, \ldots, x_n): \; x_k>0 \; \forall k$
\end{theorem}

Удивительно, но в доказательстве теоремы Фробениуса-Перрона линейная задача заменяется на нелинейную.

\begin{proof}
\begin{enumerate}
    \item $\R^n_+$~--- все координаты у элемента $x$ больше нуля: $K = \{x \in \R^n : \forall k\ x_k \geqslant 0\}$
    \item \begin{statement}
        $\forall x \in K, x \neq 0 \; Ax \in Int \; K$
    \end{statement}
    \item Рассмотрим плоскость $P: \sum x_{i_j} = 1$,
    \item $D = K \cap P$
    \item Рассмотрим отображение $F$ на $D$: $F(x)=\frac{Ax}{\sum (Ax)_k}$~--- нелинейное отображение (!!!), причём $\sum (Ax)_k \neq 0$ (см. пункт 2)
\end{enumerate}
\end{proof}